% =====================================================================
%  MODELLO: VERBALE ESTERNO  (incontro con un proponente o con i docenti)
%
%  1. Copiate questo file nella cartella della milestone in corso:
%       Candidatura/verbali/esterni/   RTB/verbali/esterni/   PB/verbali/esterni/
%     e chiamatelo VE-AAAA-MM-GG.tex (data dell'incontro), es.
%       RTB/verbali/esterni/VE-2026-11-05.tex
%     Dopo il merge la CI mette accanto il PDF: VE-2026-11-05.pdf
%     La copia firmata dal proponente salvatela accanto con un altro
%     nome: VE-2026-11-05-firmato.pdf (vedi GUIDA.md, ricetta 3)
%  2. Cambiate i campi qui sotto (dopo il simbolo % c'è un commento).
%  3. Scrivete il contenuto dopo \begin{document}.
%  Rispetto al verbale interno ci sono in più: i partecipanti esterni,
%  domande e risposte e, in fondo, il riquadro per la firma.
%  Guida completa: GUIDA.md (ricetta 2).
% =====================================================================
\documentclass[verbale=esterno]{gpt}

\titolo{Incontro conoscitivo con VIMAR}    % es. Chiarimenti sul capitolato C2
\data{2026-10-22}                  % data dell'incontro: AAAA-MM-GG
\orario{16:00}{17:00}              % ora di inizio, ora di fine
\luogo{Microsoft Teams}            % oppure: sede dell'azienda, Discord...
\redattore{mazzurana}                  % chi scrive il verbale (la sua chiave)
\verificatore{cescobolla}          % chi lo controlla (diverso dal redattore)

% Chi faceva parte del gruppo il giorno dell'incontro.
% Lasciate così, a meno che il gruppo sia cambiato.
\componenti{biasini,cescobolla,corro,guglielmo,mazzurana,spagnoli}
\assenti{}                         % chi mancava, es. \assenti{spagnoli}

% Persone esterne presenti: una riga per persona
\esterno{Mariano Sciacco}{Referente, VIMAR}
\esterno{Domenico Bevilacqua}{Referente, VIMAR}
% \esterno{Altro Nome}{Altro ruolo, Azienda}

% Chi firma il verbale per l'azienda (compare nel riquadro firma finale)
%\firma{Nome Cognome}{Azienda}

% Destinatari: ai docenti aggiungete l'azienda (\newline = a capo)
\destinatari{Prof.~Tullio Vardanega\newline Prof.~Riccardo Cardin\newline VIMAR}

% Solo DOPO l'assegnazione del capitolato (togliete il % davanti):
% \capitolato{Cn -- Nome del capitolato (Azienda)}

% ---------------------------------------------------------------------
%  REGISTRO DELLE MODIFICHE  --  la riga più recente va IN CIMA.
%  La versione del documento è quella della PRIMA riga.
%
%  \modifica{versione}{data}{autore}{verificatore}{cosa è cambiato}
%
%  Quando il Verificatore approva, aggiunge lui IN CIMA questa riga:
%  
% ---------------------------------------------------------------------
\modifica{1.0.0}{2026-10-23}{mazzurana}{cescobolla}{Verifica e approvazione}
\modifica{0.1.0}{2026-10-22}{mazzurana}{cescobolla}{Prima stesura}

\begin{document}

\section{Ordine del giorno}
\begin{enumerate}
  \item Risposte del proponente ai temi emersi dal sondaggio tra i gruppi.
  \item Domande libere (Q\&A dal vivo).
  \item Prossimi passi: redazione e firma del verbale esterno.
\end{enumerate}

\section{Svolgimento}
L'incontro è stato organizzato da Vimar S.p.A. prima dell'assegnazione
dei capitolati ed era aperto a tutti i gruppi interessati al capitolato
C1 -- AURA. Prima dell'incontro i gruppi avevano inviato le proprie
domande tramite un sondaggio: otto gruppi hanno posto 37 domande, che il
proponente ha raggruppato in 14 temi e quattro aree (architettura e
linguaggi, tecnologie e sicurezza, sviluppo e test, organizzazione).

Il Git Push Team aveva posto domande sui temi dei comandi critici e
degli argomenti da studiare subito (temi 6 e 13). Le risposte a tutti i
temi sono riportate nella sezione seguente.

\section{Domande e risposte}

\subsection{Architettura e linguaggi}

\domanda{Quale linguaggio e quale SDK vanno usati per il server MCP?}{%
  Il linguaggio è libero (Python, Java o TypeScript): il proponente non
  pone vincoli. È consigliato Python con FastMCP, l'SDK suggerito nel
  capitolato e quello su cui Vimar può offrire il supporto più diretto.
  Usare lo stesso SDK per client e server semplifica i test di
  integrazione.}

\domanda{Che ruolo deve avere il client MCP?}{%
  Lo scopo finale è che un LLM possa usare il server MCP del gruppo; i
  client commerciali (Claude, ChatGPT\ldots) includono già un client MCP
  standard. Un client proprio serve soprattutto per i test di
  integrazione; CLI, TUI ed estensione per VS Code sono opzionali. La
  divisione della gestione degli errori tra client e server (es.
  dispositivo non raggiungibile) va decisa in progettazione. Il punto
  principale resta il controllo in lettura e scrittura dei dispositivi.}

\domanda{Cosa viene fornito all'utente finale per collegare il proprio
  LLM?}{%
  Non serve sviluppare un'app né un sito: la soluzione è principalmente
  back-end e l'utente usa l'interfaccia dei servizi LLM. Si consiglia un
  server MCP locale per sviluppo e debug e, per l'utente, un server MCP
  remoto con URL pubblico e autenticazione, usabile da web, desktop e
  mobile. Basta che funzioni con un servizio AI commerciale; supportarne
  più di uno è un valore aggiunto, senza interfacce di selezione.}

\domanda{Quali risorse LLM e cloud saranno disponibili e cosa serve per
  la consegna?}{%
  Si possono usare i piani gratuiti o per studenti dei servizi
  commerciali (es. connettore personalizzato di Claude, Copilot con il
  GitHub Student Pack). Su richiesta, da discutere all'avvio, Vimar può
  dare accesso a un insieme ristretto di servizi AWS e a una chiave LLM
  per le prove. Per la consegna basta che la soluzione giri in locale con
  Docker: al collaudo Vimar la installa su AWS o su un proprio server.
  Con il GitHub Student Pack sono disponibili 100\$ di Microsoft Azure,
  in alternativa al piano gratuito di AWS.}

\subsection{Tecnologie e sicurezza}

\domanda{Come funziona l'autenticazione MyVimar in chat?}{%
  Serve un account MyVimar con accesso all'impianto; ogni utente vede
  solo i propri impianti, anche più di uno. Il server MCP associa ogni
  utente alla sua autorizzazione OAuth2 e la usa verso le API: il
  servizio AI non vede mai i token. Il flusso di login (es. link di
  autorizzazione da aprire nel browser, permessi da concedere) va
  progettato dal gruppo. I dati dell'impianto possono arrivare all'LLM
  per elaborare la risposta, i dati sensibili mai.}

\domanda{Comandi critici: chi blocca le richieste sbagliate?}{%
  Le API Vimar non hanno protezioni specifiche per l'uso con un LLM:
  leggono e scrivono come qualsiasi app. Le letture possono essere
  dirette; per le scritture (es. ``spegni tutto'') si può prevedere una
  conferma dell'utente. Gli \emph{scope} KNX IoT permettono permessi
  configurabili (sola lettura, lettura e scrittura, solo alcuni
  dispositivi). Guardrail o skill aggiuntive sono ammessi se motivati nei
  requisiti.}

\domanda{Per lo stato e le notifiche serve il polling o la
  subscription?}{%
  Per questo progetto il polling va bene, come previsto dal capitolato.
  In genere un servizio AI non riceve messaggi avviati da un webhook, ma
  si può esplorare se qualche servizio lo consente. Il caso \#21
  (``avvisami quando\ldots'') è un nice-to-have realizzabile come azione
  programmata; le subscription vanno valutate solo dopo aver completato i
  casi base.}

\domanda{Quali richieste deve saper gestire l'assistente?}{%
  Non solo le frasi degli esempi: l'LLM deve gestire varianti, sinonimi
  e richieste ambigue. I casi base e avanzati del capitolato sono un
  punto di partenza per i requisiti, non un elenco esaustivo; nuove
  interazioni sono un nice-to-have gradito. Per aumentare la variabilità
  si consiglia di far riformulare a un LLM le domande della test list e
  di provarle con modelli diversi.}

\subsection{Sviluppo e test}

\domanda{Come si usa e si condivide il kit di impianto?}{%
  Il kit si collega al Wi-Fi e si comanda via cloud con le API KNX IoT,
  senza USB né VPN. Un membro del gruppo lo tiene a casa, sempre acceso
  e collegato, e gli altri lavorano da remoto tramite le API. Vimar
  fornisce un account amministratore a cui aggiungere altri utenti
  registrati su MyVimar.}

\domanda{Esiste un ambiente simulato per sviluppare in parallelo?}{%
  Dalla specifica OpenAPI di KNX IoT si può generare in locale un mock
  server (es. con Prism); esiste anche una demo online con API di
  mocking, inclusa nella documentazione per integratori. Il mock non
  replica esattamente il concetto di impianto di Vimar: la validazione
  finale va fatta sul kit. Il mock serve per lo sviluppo quotidiano, il
  kit per i test end-to-end.}

\domanda{Cosa misurano i test end-to-end e come si gestisce la CI?}{%
  La copertura dell'80\% dei test end-to-end si misura sui requisiti di
  progetto; la test list del capitolato copre i casi base. Nei test di
  unità e di integrazione si può simulare il servizio AI, così i test
  sono ripetibili e senza costi; per questi vale la copertura del codice,
  almeno 75\% e facoltativamente oltre il 90\%. Un LLM reale non è
  deterministico, quindi è meglio tenerlo fuori dalla CI.}

\subsection{Organizzazione}

\domanda{Quali sono le priorità e cosa ci si aspetta dal PoC?}{%
  Il capitolato descrive esigenze di business: requisiti formali e
  priorità li definisce il gruppo nella propria analisi. Il PoC minimo è
  un server MCP funzionante su due o tre casi base: accesso
  all'impianto, elenco degli impianti associati, stato di una luce. Né
  gli opzionali né i nice-to-have sono obbligatori; i nice-to-have sono
  preferibili.}

\domanda{Cosa conviene studiare subito e quali difficoltà hanno avuto i
  gruppi degli anni scorsi?}{%
  Non c'è una scadenza: il lavoro parte dopo l'assegnazione, ma
  esplorare ora le tecnologie aiuta a scegliere. Conviene partire dalla
  documentazione KNX IoT per integratori, dalle primitive Vimar rispetto
  al protocollo e dallo standard MCP. Le difficoltà principali degli
  anni scorsi sono state requisiti poco chiari all'inizio,
  documentazione migliorabile e i limiti degli LLM eseguiti in locale:
  meglio chiarire presto i dubbi con Vimar e usare i servizi commerciali.}

\domanda{Come si svolgono SAL e incontri durante il progetto?}{%
  I SAL si tengono di solito su Teams, con la possibilità di incontri
  in sede Vimar. Sono previsti due incontri obbligatori in presenza
  (consegna del kit e collaudo finale) e indicativamente due demo con i
  colleghi di R\&D, al PoC e all'MVP. Vimar è disponibile per dubbi
  progettuali e scelte architetturali, ma non può revisionare i documenti
  per la parte accademica. Approfondimenti mirati vanno richiesti con una
  o due settimane di anticipo.}

\subsection{Domande libere}
Nella seconda parte dell'incontro i gruppi hanno potuto porre domande
libere tramite il Q\&A di Teams o intervenendo a voce.
% Aggiungete qui le domande fatte dal vivo, per esempio:
% \domanda{La domanda}{La risposta ricevuta.}

\section{Prossimi passi}
Il gruppo procede alla redazione del verbale esterno e all'invio da VIMAR per la firma digitale.

\section{Decisioni}
\begin{decisioni}
  \decisione{Il gruppo redige il verbale esterno dell'incontro e lo
    invia ai referenti Vimar per l'approvazione con firma digitale.}
  \decisione{Le risposte ricevute vengono usate nella valutazione dei
    capitolati in vista della scelta.}
\end{decisioni}


\section{Azioni da intraprendere}
\begin{azioni}
  \azione{Inviare il verbale ai referenti Vimar per la firma
    digitale.}{mazzurana}{2026/10/22}
  \azione{Approfondire la documentazione KNX IoT per integratori e lo
    standard MCP, per valutare il capitolato C1.}{tutti}{2026/10/22}
\end{azioni}

\end{document}
